\section{Collective Validity in Counterfactual Queries}
\label{sec:set}

While instance-level semantics define feasibility and transformation cost via $\omega$, counterfactual querying is inherently set-valued: users seek a small set of alternatives that reveal complementary actionable paths. This section introduces set-level optimization to construct informative answers, extending retrieval to structured subset selection and enabling reasoning about diversity, robustness, and fairness. It complements the challenges in Section~\ref{sec:intro} by extending efficient retrieval to set selection and enabling principled reasoning across candidates.

This requirement is captured by the principle of collective validity (P4), where answer quality depends not only on individual feasibility and cost, but also on how answers interact as a set. Simply returning the top-$k$ lowest-cost instances often yields redundant options. Instead, we define a set-level objective $\Omega$ to evaluate candidate sets, balancing cost with collective utility.

\paragraph{Robustness.}
Robust counterfactuals should remain stable under small perturbations of the query~\cite{guyomard2023generating}, ensuring that suggested actions are not overly sensitive to minor variations. Let $\mathcal N(\mathbf{x}_q)$ denote a neighborhood of the query instance. A robustness objective is,
\begingroup\small
\begin{equation}
\Omega_{\text{rob}}(S) = -\mathbb E_{\mathbf{x}_q'\sim\mathcal N(\mathbf{x}_q)} \left[ \min_{\mathbf{x}\in S} \{d(\mathbf{x}_q',\mathbf{x})\} \right].
\end{equation}
\endgroup
For each perturbed query $\mathbf{x}_q'$, the closest instance in $S$ is selected, and the expectation aggregates this distance over the neighborhood. Maximizing $\Omega_{\text{rob}}(S)$ encourages $S$ to remain representative under local variations, producing stable alternatives. Robustness supports \emph{collective validity} (P4) by ensuring the set maintains utility across a neighborhood of plausible inputs. This aligns with actionable and trustworthy counterfactual reasoning, where recommendations remain consistent and dependable under small uncertainties.

\paragraph{Diversity.}
Counterfactuals should expose distinct paths of change rather than minor variations along a single dimension~\cite{laugel2023achieving}. A diversity objective can be defined as the total pairwise dispersion:
\begingroup\small
\begin{equation}
    \Omega_{\text{div}}(S) 
    = \sum_{\mathbf{x}_i, \mathbf{x}_j \in S,\, i<j} d(\mathbf{x}_i, \mathbf{x}_j).
\end{equation}
\endgroup
Maximizing $\Omega_{\text{div}}(S)$ encourages diverse alternatives, enabling users to explore multiple recourse strategies while maintaining proximity for actionability. Diversity supports \emph{collective validity} (P4) by promoting complementary rather than redundant answers, ensuring the set spans multiple meaningful directions. This aligns with data-centric counterfactual analysis, providing informative alternative sets. This formulation corresponds to the max-sum diversification problem~\cite{borodin2017max}, with efficient approximation algorithms available.

\paragraph{Fairness.}
For auditing and decision support, the composition of the returned set is critical~\cite{kusner2017counterfactual}. Let $g(\mathbf{x})$ denote the value of a protected attribute for instance $\mathbf{x}$ (e.g., gender or ethnicity), and let $\pi_a$ denote the desired proportion of instances in the returned set that belong to group $a$ (i.e., those with $g(\mathbf{x}) = a$). A fairness objective can be defined as
\begingroup\small
\begin{equation}
\Omega_{\text{fair}}(S)
=
-
\left|
\frac{1}{|S|}\sum_{\mathbf{x}\in S}\mathbf{1}[g(\mathbf{x})=a] - \pi_a
\right|.
\end{equation}
\endgroup
Maximizing $\Omega_{\text{fair}}(S)$ encourages the selected set to match desired group proportions, reducing bias in the returned explanations. Fairness supports \emph{collective validity} (P4) by coupling decisions across candidates and introducing global dependencies. This aligns with responsible, human-centered counterfactual reasoning, where outputs should be feasible, informative, and equitable. In practice, enforcing fairness increases combinatorial complexity and must be balanced with objectives such as cost and diversity.
